\subsection{Tensor reconstruction and projection of sextic centrifugal distortion constants}

The sextic case can be treated within the same tensor-based framework,
although the corresponding tensor space is substantially richer than in the
quartic case.  Let
\[
\mathbf D^{(6)} =
\bigl(
H_J,H_{JK},H_{KJ},H_K,h_1,h_2,h_3
\bigr)^T
\]
denote the sextic Watson constants in a given reduction and principal-axis
representation.  If the starting point is a set of spectroscopic constants, the
corresponding reduced sextic tensor representation
$\boldsymbol\phi^{(6)}$ is reconstructed through
\begin{equation}
\boldsymbol\phi^{(6)}=
\mathbf R_6^{+}(A,B,C)\,\mathbf D^{(6)},
\label{eq:sextic_pseudoinverse}
\end{equation}
where $\mathbf R_6^{+}$ is the Moore--Penrose pseudoinverse of the sextic
projection matrix.  As in the quartic case, the pseudoinverse enters only in
this inverse reconstruction step.

For the practical analysis of representation changes it is convenient to work
in Watson's $S$ reduction and to decompose the resulting seven-component sextic
vector as
\begin{equation}
\mathbf H_S = \mathbf B_{\mathrm{rep}}\,\boldsymbol\sigma + \mathbf r,
\label{eq:sextic_5plus2}
\end{equation}
where $\mathbf B_{\mathrm{rep}}$ is a representation-dependent $7\times 5$
matrix, $\boldsymbol\sigma$ is a five-component coordinate vector on the
canonical sextic subspace, and $\mathbf r$ is the residual in the complementary
two-dimensional sector.  For internally consistent sextic parameter sets one
has
\begin{equation}
\mathbf r=\mathbf 0
\end{equation}
within numerical precision, so that the residual serves as a direct
diagnostic of the distance from the canonical physical subspace.

The five-dimensional sextic sector used in the present work is not an
arbitrary numerical reduction.  The even degree-six rotational polynomial space
admits the harmonic decomposition
\begin{equation}
\mathcal P_6^{(\mathrm{even})}
=
r^6\mathcal H_0
\oplus
r^4\mathcal H_2^{(\mathrm{even})}
\oplus
r^2\mathcal H_4^{(\mathrm{even})}
\oplus
\mathcal H_6^{(\mathrm{even})},
\label{eq:sextic_harmonic_decomp}
\end{equation}
with dimensions
\[
10 = 1 + 2 + 3 + 4.
\]
Here $r^2=J_a^2+J_b^2+J_c^2$, $\mathcal H_0$ is the isotropic scalar sector,
and $\mathcal H_\ell^{(\mathrm{even})}$ denotes the even-parity harmonic sector
of degree $\ell$.

For every sextic even polynomial $p$, the harmonic decomposition
\[
p = h_6 + r^2 h_4 + r^4 h_2 + r^6 h_0
\]
is unique, with
\[
h_6\in\mathcal H_6^{(\mathrm{even})},\qquad
h_4\in\mathcal H_4^{(\mathrm{even})},\qquad
h_2\in\mathcal H_2^{(\mathrm{even})},\qquad
h_0\in\mathcal H_0.
\]
If $\Delta$ denotes the Laplacian in $(J_a,J_b,J_c)$, these components are
extracted canonically through
\begin{align}
h_0 &= \frac{\Delta^3 p}{5040},\\
h_2 &= \frac{\Delta^2 p}{504}-\frac{r^2\Delta^3 p}{3024},\\
h_4 &= \frac{\Delta p}{22}-\frac{r^2\Delta^2 p}{308}
+\frac{r^4\Delta^3 p}{6160},\\
h_6 &= p-\frac{r^2\Delta p}{22}
+\frac{r^4\Delta^2 p}{792}
-\frac{r^6\Delta^3 p}{33264}.
\end{align}

The canonical five-dimensional sextic subspace employed in the present
implementation projects isomorphically onto
\[
\mathcal H_6^{(\mathrm{even})}\oplus r^6\mathcal H_0.
\]
Consequently, the physical sextic content is equivalently described by
\begin{equation}
\boldsymbol\sigma \;\longleftrightarrow\;
(\sigma_0,\sigma_1,\sigma_2,\sigma_3,\sigma_4),
\label{eq:sextic_1plus4}
\end{equation}
that is, by one isotropic amplitude and four amplitudes spanning the even
harmonic degree-six sector.  In this sense the canonical sextic subspace admits
a natural $1+4$ interpretation.  The lower-order sectors
$r^2\mathcal H_4^{(\mathrm{even})}$ and
$r^4\mathcal H_2^{(\mathrm{even})}$ do not carry independent physical
coordinates; rather, they are determined by the graph relation defining the
canonical sextic subspace.  The operational decomposition
Eq.~(\ref{eq:sextic_5plus2}) therefore remains the most convenient form for the
representation transform, whereas Eq.~(\ref{eq:sextic_1plus4}) provides its
underlying structural interpretation.  Equivalently, the physical sextic
content may be identified with the pair $(h_0,h_6)$, while the lower harmonic
pieces $(h_2,h_4)$ are induced from it by the fixed graph relation defining the
canonical sextic sector.  Here $h_6$ is not merely a four-dimensional block:
it is precisely the $D_2$-even sector of the irreducible $\ell=6$
representation of $SO(3)$.  The canonical sextic sector therefore admits a
more geometric interpretation as one isotropic scalar amplitude plus one even
rank-six irreducible harmonic component.  Choosing any real tesseral basis
\[
\{Y_6^{(0)},Y_6^{(2c)},Y_6^{(4c)},Y_6^{(6c)}\}
\]
for this $D_2$-even $\ell=6$ sector, one may write
\[
h_6=c_0Y_6^{(0)}+c_2Y_6^{(2c)}+c_4Y_6^{(4c)}+c_6Y_6^{(6c)}.
\]
The operational coordinates $\boldsymbol\sigma$ are then linearly equivalent to
the harmonic amplitude vector
\[
(h_0,c_0,c_2,c_4,c_6),
\]
so that the present sextic formalism may be interpreted not only as a
canonical $5+2$ decomposition, but also as a basis-independent
scalar-plus-rank-six irreducible parametrization of the physical sextic
content.

For CeDiTT, however, one further refinement is especially natural.  The
representation changes implemented in the program are generated by cyclic
relabelings of the principal axes,
\[
I^r \rightarrow II^r \rightarrow III^r \rightarrow I^r,
\]
that is, by the cyclic group $C_3$.  Under this action the
four-dimensional space $\mathcal H_6^{(\mathrm{even})}$ is not irreducible:
it decomposes as
\[
\mathcal H_6^{(\mathrm{even})} \cong 1 \oplus 1 \oplus \rho_2,
\]
where the first two summands are one-dimensional invariant directions and
$\rho_2$ is the real two-dimensional irreducible representation of $C_3$.
Accordingly, the physical sextic sector admits a representation-adapted
decomposition
\begin{equation}
\mathcal H_0 \oplus \mathcal H_6^{(\mathrm{even})}
\cong
1 \oplus 1 \oplus 1 \oplus 2.
\label{eq:sextic_rep_adapted_3plus2}
\end{equation}
This provides the natural origin of the effective $3+2$ structure in the
sextic problem: three scalar coordinates remain invariant under cyclic
representation changes, whereas the remaining two coordinates form a doublet
that mixes by a planar rotation in the two-dimensional irrep $\rho_2$.  The
coordinates $\boldsymbol\sigma$ used in the present implementation provide one
convenient basis for this canonical representation-adapted structure.

Once the sextic tensor coefficients have been reconstructed, the forward
projection is most conveniently written in terms of the Cartesian sextic
coefficients in the principal-axis frame,
\[
\Phi_{aaa},\Phi_{aab},\Phi_{aac},\Phi_{abb},\Phi_{abc},\Phi_{acc},
\Phi_{bbb},\Phi_{bbc},\Phi_{bcc},\Phi_{ccc},
\]
which are the coefficients of the sextic rotational polynomial in the
molecule-fixed principal-axis system.  These coefficients are first transformed
to cylindrical sextic quantities,
\begin{align}
\Phi_{600} &= \frac{5}{16}(\Phi_{aaa}+\Phi_{bbb})
+\frac{1}{8}(\Phi_{aab}+\Phi_{abb}),\\
\Phi_{420} &= \frac{3}{4}(\Phi_{aac}+\Phi_{bbc})
+\frac{1}{4}\Phi_{abc}-3\Phi_{600},\\
\Phi_{240} &= \Phi_{acc}+\Phi_{bcc}-2\Phi_{420}-3\Phi_{600},\\
\Phi_{060} &= \Phi_{ccc}-\Phi_{240}-\Phi_{420}-\Phi_{600},\\
\Phi_{402} &= \frac{15}{64}(\Phi_{aaa}-\Phi_{bbb})
+\frac{1}{32}(\Phi_{aab}-\Phi_{abb}),\\
\Phi_{222} &= \frac{1}{2}(\Phi_{aac}-\Phi_{bbc})-2\Phi_{402},\\
\Phi_{042} &= \frac{1}{2}(\Phi_{acc}-\Phi_{bcc})
-\Phi_{222}-\Phi_{402},\\
\Phi_{204} &= \frac{3}{32}(\Phi_{aaa}+\Phi_{bbb})
-\frac{1}{16}(\Phi_{aab}+\Phi_{abb}),\\
\Phi_{024} &= \frac{1}{8}(\Phi_{aac}+\Phi_{bbc}-\Phi_{abc})
-\Phi_{204},\\
\Phi_{006} &= \frac{1}{64}(\Phi_{aaa}-\Phi_{bbb})
-\frac{1}{32}(\Phi_{aab}-\Phi_{abb}).
\end{align}

The quartic cylindrical quantities enter the sextic projection through
\begin{align}
B_{200} &= \frac{1}{2}(A+B)-4T_{004},\\
B_{020} &= C-B_{200}+6T_{004},\\
B_{002} &= \frac{1}{4}(A-B),
\end{align}
reflecting the perturbative structure by which sextic centrifugal distortion
contains quadratic contributions from the quartic rotational tensor.

For the $A$ reduction, the sextic constants are
\begin{align}
\Phi_J &= \Phi_{600}+2\Phi_{204},\\
\Phi_{JK} &= \Phi_{420}-12\Phi_{204}+2\Phi_{024}
+16\Sigma\Phi_{006}
+8\frac{T_{022}T_{004}}{B_{002}},\\
\Phi_{KJ} &= \Phi_{240}+\frac{10}{3}\Phi_{420}
-30\Phi_{204}-\frac{10}{3}\Phi_{JK},\\
\Phi_K &= \Phi_{060}-\frac{7}{3}\Phi_{420}
+28\Phi_{204}+\frac{7}{3}\Phi_{JK},\\
\phi_J &= \Phi_{402}+\Phi_{006},\\
\phi_K &= \Phi_{042}+\frac{4}{3}\Sigma\Phi_{024}
+\left(\frac{32}{3}\Sigma^2+9\right)\Phi_{006}
+4\frac{T_{040}+\frac{1}{3}\Sigma T_{022}
-2(\Sigma^2-2)T_{004}}{B_{002}}T_{004},\\
\phi_{JK} &= \Phi_{222}+4\Sigma\Phi_{204}-10\Phi_{006}
+2\frac{T_{220}-2\Sigma T_{202}-4T_{004}}{B_{002}}T_{004},
\end{align}
with
\begin{equation}
\Sigma=\frac{2A-B-C}{B-C}.
\end{equation}

For the $S$ reduction one defines
\begin{align}
\rho &= \frac{1}{4}T_{022}\Sigma_1,\\
\mu &= \frac{\Sigma\Phi_{042}-\frac{9}{8}\Phi_{024}
+\bigl(-2\Sigma T_{040}+(\Sigma^2+3)T_{022}-5\Sigma T_{004}\bigr)
T_{022}/B_{020}}
{2\Sigma^2+\frac{27}{16}},\\
\nu &= \frac{3}{16}\mu\Sigma_1+\frac{1}{8}\Phi_{024}\Sigma_1
+\frac{T_{004}T_{022}}{B_{020}},\\
\lambda &= 5\nu\Sigma_1+\frac{1}{2}\Phi_{222}\Sigma_1
+\frac{-\frac{1}{2}T_{220}\Sigma_1+T_{202}-T_{022}\Sigma_1^2
-2T_{004}\Sigma_1}{B_{020}}T_{022},
\end{align}
with
\begin{equation}
\Sigma_1=\frac{B-C}{2A-B-C},
\end{equation}
and finally
\begin{align}
H_J &= \Phi_{600}-\lambda,\\
H_{JK} &= \Phi_{420}+6\lambda-3\mu,\\
H_{KJ} &= \Phi_{240}-5\lambda+10\mu,\\
H_K &= \Phi_{060}-7\mu,\\
h_1 &= \Phi_{402}-\nu,\\
h_2 &= \Phi_{204}+\frac{1}{2}\lambda,\\
h_3 &= \Phi_{006}+\nu.
\end{align}

These equations provide the explicit forward projection from sextic tensor
coefficients to Watson sextic constants.  Thus, exactly as in the quartic case,
the pseudoinverse is required only when one reconstructs the tensor from
spectroscopic constants.  Once the sextic tensor coefficients are known, the
projection to Watson constants is fully explicit.  Conversely, when only Watson
constants are available, Eq.~(\ref{eq:sextic_5plus2}) provides the operational
decomposition used in the representation transform, while
Eq.~(\ref{eq:sextic_1plus4}) identifies the same physical sector with one
isotropic scalar amplitude and four amplitudes in the even sextic harmonic
sector.  In the representation-adapted viewpoint of
Eq.~(\ref{eq:sextic_rep_adapted_3plus2}), the same five-dimensional content is
resolved into three scalar invariants and one two-dimensional doublet under the
cyclic relabeling of the principal axes.
sector.
